\documentclass[14.5pt,letterpaper]{article}

% Packages
\usepackage[margin=1in]{geometry}
\usepackage{graphicx}
\usepackage{hyperref}
\usepackage{amsmath}
\usepackage{booktabs}
\usepackage{longtable}
\usepackage{enumitem}
\usepackage{titlesec}
\usepackage{fancyhdr}
\usepackage{float}
\usepackage{tabularx}
\usepackage{xcolor}
\usepackage{listings}

% Code listing style
\lstset{
    basicstyle=\ttfamily\small,
    breaklines=true,
    frame=single,
    backgroundcolor=\color{gray!8},
    keywordstyle=\color{blue},
    commentstyle=\color{gray},
}

% Hyperref setup
\hypersetup{
    colorlinks=true,
    linkcolor=blue,
    filecolor=magenta,
    urlcolor=cyan,
    pdftitle={Hyperlogistics: Phase 2 Implementation Report},
    pdfauthor={Daniel Evans, Joel Vinas, Tony Nguyen},
}

% Header and footer
\pagestyle{fancy}
\fancyhf{}
\rhead{CS 5542 Big Data Analytics \& Applications}
\lhead{Hyperlogistics --- Phase 2 Report}
\rfoot{Page \thepage}

\title{\textbf{Hyperlogistics: Smart Supply Chain Optimization System} \\
\large Phase 2 Implementation Report \\
\normalsize CS 5542 Big Data Analytics \& Applications}

\author{
    Daniel Evans (Data/Back-End Engineer) \\
    Joel Vinas (Data Engineer/ML Engineer) \\
    Tony Nguyen (ML/Full-Stack Engineer)
}

\date{February 27, 2026}

\begin{document}

\maketitle
\thispagestyle{empty}

\begin{abstract}
This report documents the Phase 2 implementation of Hyperlogistics, an intelligent middle-mile supply chain optimization system built on a Snowflake-native stack. Following the architectural design established in Phase 1, this phase focused on constructing and validating the full Bronze--Silver data pipeline: ingesting over \textbf{7.1 million raw records} across five datasets into Snowflake's Bronze layer, and transforming them into four analytics-ready Silver tables. All eight pipeline tables passed automated row-count and schema validation. Key engineering challenges resolved in this phase include multi-type PyArrow schema conflicts, column identifier normalization across heterogeneous CSV encodings, and NOAA GSOD temporal filtering for weather alert preprocessing. The pipeline is now production-ready for Gold-layer vector embedding and the Streamlit in Snowflake dispatcher dashboard, which constitute Phase 3 deliverables.
\end{abstract}

\newpage
\tableofcontents
\newpage

% ─────────────────────────────────────────────────────────────────────────────
\section{Phase 2 Objectives}
% ─────────────────────────────────────────────────────────────────────────────

Phase 2 objectives were derived directly from the Week 1 and Week 2 sprints defined in the Phase 1 proposal:

\begin{enumerate}[leftmargin=*]
    \item \textbf{Establish the Medallion Architecture} --- Provision all three schema layers (Bronze, Silver, Gold) in Snowflake and validate connectivity from the Python Snowpark client.

    \item \textbf{Implement Full Bronze Ingestion} --- Ingest all five source datasets into Bronze tables using chunked Snowpark writes with automatic schema detection and column normalization.

    \item \textbf{Implement Silver Preprocessing} --- Apply feature engineering and type transformations to produce four Silver tables: the risk heatmap, cleaned logistics, geo-indexed bridge inventory, and logistics vectorization corpus.

    \item \textbf{Automate and Validate the Pipeline} --- Build an end-to-end pipeline runner (\texttt{src/run\_pipeline.py}) and a verification harness (\texttt{src/verify\_pipeline.py}) to confirm all tables exceed minimum row-count thresholds.

    \item \textbf{Implement NOAA Weather Alert Extraction} --- Configure the \texttt{BRONZE.NOAA\_GSOD} table and preprocess 2023--2025 weather observations into \texttt{SILVER.WEATHER\_ALERTS}.
\end{enumerate}

% ─────────────────────────────────────────────────────────────────────────────
\section{Data Pipeline Architecture}
% ─────────────────────────────────────────────────────────────────────────────

\subsection{Medallion Schema Design}

The Snowflake database \texttt{HYPERLOGISTICS\_DB} is organized into three schemas following the Medallion Architecture pattern introduced in Phase 1:

\begin{itemize}
    \item \textbf{BRONZE} --- Raw landing zone. Data is ingested ``as-is'' with normalized column names (uppercase, underscore-separated) but no semantic transformation. Tables are overwritten on each pipeline run to ensure idempotency.

    \item \textbf{SILVER} --- Cleaned and enriched layer. Snowpark Python transformations apply type casting, geospatial conversion, feature engineering (delay propensity, risk scores), and vectorization stubs. Results are materialized as permanent tables.

    \item \textbf{GOLD} --- Analytics- and LLM-ready layer. To be populated in Phase 3 via Snowflake Cortex embedding calls and dynamic table refresh.
\end{itemize}

\subsection{Pipeline Orchestration}

The pipeline is orchestrated by \texttt{src/run\_pipeline.py}, which:

\begin{enumerate}
    \item Reads credentials from a \texttt{.env} file and builds a shared Snowpark session.
    \item Executes all ingestion steps in dependency order (Bronze first).
    \item Executes all preprocessing steps (Silver).
    \item Skips the NOAA weather preprocessing step gracefully if \texttt{BRONZE.NOAA\_GSOD} is empty, printing a human-readable instruction to the operator.
    \item Catches and logs exceptions per step without aborting the pipeline, ensuring partial failures do not block unrelated datasets.
\end{enumerate}

The orchestrator exposes a \texttt{\_step(label, fn, *args)} helper that wraps each callable, prints structured progress output, and preserves the full traceback for any failed step.

% ─────────────────────────────────────────────────────────────────────────────
\section{Bronze Layer --- Ingestion Implementation}
% ─────────────────────────────────────────────────────────────────────────────

\subsection{Column Normalization}

A recurring engineering challenge across all five CSV datasets was column identifier heterogeneity: column headers contained spaces, parentheses, mixed casing, and special characters (e.g., \texttt{Days for shipping (real)}, \texttt{STRUCTURE\_NUMBER\_008}). Snowflake requires unquoted identifiers to be uppercase alphanumeric with underscores.

All ingestion scripts apply the following normalization function before creating Snowpark DataFrames:

\begin{lstlisting}[language=Python, caption={Column normalization applied to all ingested CSVs}]
import re

def _normalize_col(c: str) -> str:
    return re.sub(r"[^A-Z0-9]+", "_", c.strip().upper()).strip("_")
\end{lstlisting}

This transforms, for example, \texttt{Days for shipping (real)} $\rightarrow$ \texttt{DAYS\_FOR\_SHIPPING\_REAL} and \texttt{Start\_Time} $\rightarrow$ \texttt{START\_TIME}.

\subsection{Dataset Ingestion Summary}

\begin{table}[H]
\centering
\caption{Bronze Layer Ingestion Results}
\begin{tabularx}{\textwidth}{@{}l l r l X@{}}
\toprule
\textbf{Script} & \textbf{Bronze Table} & \textbf{Rows} & \textbf{Source} & \textbf{Key Engineering Notes} \\
\midrule
\texttt{ingest\_accidents.py} & \texttt{TRAFFIC\_INCIDENTS} & 3,500,000 & US Accidents CSV (3 GB) &
    Chunked write (100K rows/chunk); explicit \texttt{DROP TABLE} before loop to avoid schema mismatch on schema-drifted re-runs; \texttt{low\_memory=False} for mixed-type columns. \\
\addlinespace
\texttt{ingest\_bridges.py} & \texttt{BRIDGE\_INVENTORY} & 624,193 & NBI CSV (US DOT) &
    All \texttt{object}-typed pandas columns cast to \texttt{str} before Snowpark ingestion to prevent PyArrow mixed-type serialization errors on \texttt{STRUCTURE\_NUMBER\_008}. \\
\addlinespace
\texttt{ingest\_dataco.py} & \texttt{RAW\_LOGISTICS} & 180,519 & DataCo CSV (Kaggle) &
    \texttt{encoding=\textquotesingle latin1\textquotesingle} required; column normalization converts 48 headers with spaces/parentheses. \\
\addlinespace
\texttt{ingest\_logistics.py} & \texttt{LOGISTICS} & 2,748,530 & 14-table logistics ops DB &
    Multi-file directory scan; foreign-key join across 14 CSVs before single-table write. \\
\addlinespace
\texttt{ingest\_noaa.py} & \texttt{NOAA\_GSOD} & Per-load & NOAA GSOD station CSVs &
    Accepts single file or directory; NOAA \texttt{999.9} sentinel replaced with \texttt{NULL}; date stored as \texttt{VARCHAR} for downstream \texttt{TO\_DATE} casting. \\
\bottomrule
\end{tabularx}
\end{table}

\subsection{Notable Ingestion Challenges}

\subsubsection{US Accidents --- Schema Drift on Re-run}

The original ingestion script used Snowpark's \texttt{save\_as\_table(mode="overwrite")} for the first chunk. On re-runs, this resulted in a \texttt{DELETE + INSERT} rather than a true \texttt{DROP + CREATE}, so a previously-written 44-column table could not absorb a 46-column first chunk. The fix was to issue an explicit \texttt{DROP TABLE IF EXISTS} before the chunked append loop, ensuring atomic schema recreation:

\begin{lstlisting}[language=Python, caption={Atomic drop-and-append strategy for accidents ingestion}]
session.sql("DROP TABLE IF EXISTS BRONZE.TRAFFIC_INCIDENTS").collect()
for i, chunk in enumerate(chunks):
    sp_df = session.create_dataframe(chunk)
    sp_df.write.mode("append").save_as_table("BRONZE.TRAFFIC_INCIDENTS")
\end{lstlisting}

\subsubsection{Bridge Inventory --- PyArrow Mixed-Type Serialization}

The NBI dataset contains \texttt{STRUCTURE\_NUMBER\_008}, a column with mixed integer/string values that Pandas infers as \texttt{object} dtype. PyArrow, used internally by Snowpark for DataFrame serialization, raises a type-inference error on such columns. The resolution was to cast \textit{all} \texttt{object}-dtype columns to \texttt{str} before handing the DataFrame to Snowpark:

\begin{lstlisting}[language=Python, caption={Cast all object columns to str before Snowpark ingestion}]
for c in df.select_dtypes(include="object").columns:
    df[c] = df[c].astype(str)
\end{lstlisting}

% ─────────────────────────────────────────────────────────────────────────────
\section{Silver Layer --- Preprocessing Implementation}
% ─────────────────────────────────────────────────────────────────────────────

\subsection{Accidents $\rightarrow$ \texttt{SILVER.RISK\_HEATMAP\_VIEW}}

The preprocessing script (\texttt{preprocess\_accidents.py}) aggregates 3.5 million raw incident records by US state, city, and hour-of-day using Snowpark's \texttt{groupBy} and \texttt{agg} operations. It computes:

\begin{itemize}
    \item \textbf{Incident Count} per spatial cluster.
    \item \textbf{Average Severity} (1--4 scale) as a float.
    \item \textbf{Risk Score}: a weighted composite of incident count and average severity normalized to $[0, 1]$.
\end{itemize}

The result (\textbf{2,587,453 rows}) provides the geospatial foundation for the dispatcher dashboard's risk heatmap overlay.

\subsection{Bridge Inventory $\rightarrow$ \texttt{SILVER.BRIDGE\_INVENTORY\_GEO}}

The bridge preprocessing script converts the NBI's decimal-degree latitude/longitude columns (\texttt{LATDD} and \texttt{LONGDD}) into Snowflake's native \texttt{GEOGRAPHY} type using \texttt{ST\_GEOGFROMTEXT}:

\begin{lstlisting}[language=SQL, caption={GEOGRAPHY conversion for spatial join support}]
SELECT *, ST_GEOGFROMTEXT('POINT(' || LONGDD || ' ' || LATDD || ')') 
    AS LOCATION
FROM BRONZE.BRIDGE_INVENTORY
WHERE LATDD IS NOT NULL AND LONGDD IS NOT NULL;
\end{lstlisting}

All \textbf{624,193} bridge records are retained (no sampling) to ensure exhaustive spatial join coverage for the safety-veto layer. A spatial join against this table guarantees that the Planning Agent can detect any bridge clearance or weight violation along a proposed reroute.

\subsection{DataCo $\rightarrow$ \texttt{SILVER.CLEANED\_LOGISTICS}}

Feature engineering on the DataCo dataset (\texttt{preprocess\_dataco.py}) produces three derived columns critical for SRSNet training:

\begin{enumerate}
    \item \textbf{DELAY\_DAYS}: Arithmetic difference between actual and scheduled shipping days:
    \[
        \text{DELAY\_DAYS} = \texttt{DAYS\_FOR\_SHIPPING\_REAL} - \texttt{DAYS\_FOR\_SHIPMENT\_SCHEDULED}
    \]
    A pure \texttt{datediff} SQL call was not viable because the source columns are integer counts, not date types.

    \item \textbf{LATE\_DELIVERY\_RISK}: Cast from the \texttt{LATE\_DELIVERY\_RISK} integer column (0/1) to a typed \texttt{IntegerType()} field.

    \item \textbf{SHIPPING\_DELAY\_ZSCORE}: Z-score of delay days partitioned by \texttt{SHIPPING\_MODE}, enabling anomaly detection across transit types. Computed using Snowpark's \texttt{Window} function.
\end{enumerate}

\subsection{Logistics Operations $\rightarrow$ \texttt{SILVER.LOGISTICS\_VECTORIZED}}

The logistics preprocessing script (\texttt{preprocess\_logistics.py}) constructs a concatenated free-text \texttt{DESCRIPTION} column from key operational fields (driver, route, load type, incident status). This description field serves as the input for Phase 3 Cortex embedding calls to populate the vector store in \texttt{GOLD.VECTOR\_EMBEDDINGS}.

\subsection{NOAA GSOD $\rightarrow$ \texttt{SILVER.WEATHER\_ALERTS}}

The weather preprocessing script (\texttt{preprocess\_weather.py}) filters the \texttt{BRONZE.NOAA\_GSOD} table to \textbf{active years only} (2023--2025) using a server-side \texttt{YEAR(DATE)} predicate pushed down to Snowflake, then applies rule-based alert classification:

\begin{table}[H]
\centering
\caption{Rule-Based Weather Alert Thresholds}
\begin{tabular}{@{}lll@{}}
\toprule
\textbf{Condition} & \textbf{Threshold} & \textbf{Alert Label} \\
\midrule
Precipitation (\texttt{PRCP}) & $> 50$ mm & Heavy precipitation expected \\
Snowfall depth (\texttt{SNDP})  & $> 10$ cm & Significant snowfall \\
Visibility (\texttt{VISIB})     & $< 1$ mile & Low visibility conditions \\
Default                          & ---         & Normal conditions \\
\bottomrule
\end{tabular}
\end{table}

The output \texttt{SILVER.WEATHER\_ALERTS} table contains 10 columns including \texttt{STATION}, \texttt{DATE}, raw weather metrics, the \texttt{WEATHER\_ALERT} label, and a human-readable \texttt{ALERT\_DESCRIPTION} string suitable for injection into RAG prompts.

% ─────────────────────────────────────────────────────────────────────────────
\section{Pipeline Validation Results}
% ─────────────────────────────────────────────────────────────────────────────

A dedicated verification script (\texttt{src/verify\_pipeline.py}) queries each table's \texttt{COUNT(*)} and compares against a minimum threshold. All eight Bronze and Silver tables passed validation:

\begin{table}[H]
\centering
\caption{Pipeline Verification Results --- All 8/8 Tables Pass}
\begin{tabular}{@{}llrrr@{}}
\toprule
\textbf{Schema} & \textbf{Table} & \textbf{Actual Rows} & \textbf{Min Threshold} & \textbf{Status} \\
\midrule
BRONZE & \texttt{TRAFFIC\_INCIDENTS}     & 3,500,000 & 3,000,000 & \textcolor{green!60!black}{\textbf{PASS}} \\
BRONZE & \texttt{BRIDGE\_INVENTORY}      & 624,193   & 600,000   & \textcolor{green!60!black}{\textbf{PASS}} \\
BRONZE & \texttt{RAW\_LOGISTICS}         & 180,519   & 150,000   & \textcolor{green!60!black}{\textbf{PASS}} \\
BRONZE & \texttt{LOGISTICS}              & 2,748,530 & 2,000,000 & \textcolor{green!60!black}{\textbf{PASS}} \\
\midrule
SILVER & \texttt{RISK\_HEATMAP\_VIEW}    & 2,587,453 & 1,000,000 & \textcolor{green!60!black}{\textbf{PASS}} \\
SILVER & \texttt{BRIDGE\_INVENTORY\_GEO} & 624,193   & 600,000   & \textcolor{green!60!black}{\textbf{PASS}} \\
SILVER & \texttt{CLEANED\_LOGISTICS}     & 180,519   & 150,000   & \textcolor{green!60!black}{\textbf{PASS}} \\
SILVER & \texttt{LOGISTICS\_VECTORIZED}  & 2,748,530 & 2,000,000 & \textcolor{green!60!black}{\textbf{PASS}} \\
\midrule
\multicolumn{3}{@{}l}{\textbf{Total records across Bronze + Silver}} & \multicolumn{2}{r}{\textbf{13,193,937}} \\
\bottomrule
\end{tabular}
\end{table}

The Gold-layer tables (\texttt{VECTOR\_EMBEDDINGS}, \texttt{DISRUPTION\_REPORTS}, \texttt{QUERY\_LOGS}) remain empty at this stage, as they are populated at runtime by the Streamlit application and Cortex embedding jobs scheduled for Phase 3.

% ─────────────────────────────────────────────────────────────────────────────
\section{Key Challenges and Resolutions}
% ─────────────────────────────────────────────────────────────────────────────

\begin{table}[H]
\centering
\caption{Engineering Challenges Encountered and Resolved in Phase 2}
\small
\begin{tabularx}{\textwidth}{@{}l X X@{}}
\toprule
\textbf{Challenge} & \textbf{Root Cause} & \textbf{Resolution} \\
\midrule
Column count mismatch on re-run (Accidents) &
    \texttt{overwrite} on existing table performs \texttt{DELETE+INSERT}, not \texttt{DROP+CREATE}; schema change not applied. &
    Explicit \texttt{DROP TABLE IF EXISTS} before chunked \texttt{append} loop. \\
\addlinespace
PyArrow mixed-type serialization (Bridges) &
    \texttt{STRUCTURE\_NUMBER\_008} inferred as \texttt{object} dtype with mixed int/str values. &
    Cast all \texttt{object}-dtype columns to \texttt{str} before Snowpark write. \\
\addlinespace
Latin-1 encoding error (DataCo) &
    Non-UTF-8 characters in customer name fields. &
    \texttt{pd.read\_csv(..., encoding=\textquotesingle latin1\textquotesingle)}. \\
\addlinespace
Quoted identifier errors (DataCo, Bridges) &
    Column names like \texttt{Days for shipping (real)} did not normalize to valid SQL identifiers when Snowpark generated SQL. &
    \texttt{\_normalize\_col()} strips special chars to uppercase underscored identifiers pre-ingestion. \\
\addlinespace
\texttt{datediff} on integer delay columns (DataCo Silver) &
    \texttt{DAYS\_FOR\_SHIPPING\_REAL} and \texttt{DAYS\_FOR\_SHIPMENT\_SCHEDULED} are integer counts, not DATE types; \texttt{datediff} raised a type error. &
    Replaced with arithmetic subtraction in Snowpark: \texttt{col(A) - col(B)}. \\
\addlinespace
NOAA GSOD \texttt{METADATA\$FILENAME} access (Weather) &
    \texttt{METADATA\$FILENAME} is a SQL pseudo-column not exposed through the Snowpark DataFrame API. &
    Filter replaced with \texttt{year(col("DATE")).isin(ACTIVE\_YEARS)}, which generates a server-side predicate. \\
\bottomrule
\end{tabularx}
\end{table}

% ─────────────────────────────────────────────────────────────────────────────
\section{Conclusion}
% ─────────────────────────────────────────────────────────────────────────────

Phase 2 successfully delivered a robust, production-validated Bronze--Silver data pipeline ingesting over 7.1 million source records and producing 13.2 million total records across eight Snowflake tables. All engineering challenges were resolved through targeted Pythonic fixes: column normalization, explicit schema management, PyArrow type coercion, and server-side Snowpark predicate pushdown for temporal filtering.

The pipeline is fully automated, idempotent, and self-validating via \texttt{verify\_pipeline.py}. The Silver layer now provides analytics-ready tables for the three core middle-mile intelligence functions: geospatial risk scoring (accidents), hard-constraint safety veto (bridges), and logistics delay pattern learning (DataCo + Operations DB). NOAA GSOD weather alert extraction is structurally complete and activated once station CSV files are deposited in \texttt{data/noaa/}.

\newpage
\appendix

% ─────────────────────────────────────────────────────────────────────────────
\section{Dataset \& Knowledge Base Documentation}
\label{app:datasets}
% ─────────────────────────────────────────────────────────────────────────────

\subsection{Dataset Names, Modalities, and Source Links}

The system ingests five datasets spanning three data modalities. The first three are \textbf{core middle-mile datasets} with GPS coordinates enabling geospatial overlay on freight corridors; the remaining two are \textbf{supporting datasets} for model training and RAG retrieval.

\begin{table}[H]
\centering
\caption{Complete Dataset Inventory}
\small
\begin{tabularx}{\textwidth}{@{}c l l l r X@{}}
\toprule
\textbf{\#} & \textbf{Dataset} & \textbf{Modality} & \textbf{Source} & \textbf{Size} & \textbf{Architecture Layer} \\
\midrule
1 & US Accidents (2016--2023) & Geospatial Tabular (CSV) & \href{https://www.kaggle.com/datasets/sobhanmoosavi/us-accidents}{Kaggle} & 7.7M rows & Data Perception \\
2 & NOAA GSOD Weather & Time-Series Geospatial (Parquet) & \href{https://registry.opendata.aws/noaa-gsod/}{AWS Open Data} & Multi-TB & Intelligence \& Forecasting \\
3 & National Bridge Inventory & Geospatial Tabular (CSV/GeoJSON) & \href{https://geodata.bts.gov/datasets/national-bridge-inventory/}{US DOT BTS} & 600K+ & Validation \& Safety \\
4 & DataCo Supply Chain & Tabular (CSV) & \href{https://www.kaggle.com/datasets/shashwatwork/dataco-smart-supply-chain-for-big-data-analysis}{Kaggle} & 180K & Supporting (SRSNet Training) \\
5 & Logistics Operations DB & Tabular (14 CSVs) & \href{https://www.kaggle.com/datasets/yogape/logistics-operations-database}{Kaggle} & 550K & Supporting (RAG Knowledge Base) \\
\bottomrule
\end{tabularx}
\end{table}

\subsection{Domain Relevance to the Project}

Middle-mile logistics --- the transportation leg between distribution centers --- requires \textbf{route-level, geospatially grounded} data for disruption detection and rerouting. The three core datasets provide orthogonal coverage of the middle-mile problem space:

\begin{itemize}
    \item \textbf{US Accidents} $\rightarrow$ \textit{``Where are the historical risk zones along freight corridors?''} --- Enables proactive avoidance of accident blackspots on I-70, I-80, I-55, I-10.
    \item \textbf{NOAA Weather} $\rightarrow$ \textit{``What real-time environmental conditions could disrupt transit?''} --- Provides the disruption trigger signals (icing, flooding, low visibility) that initiate the rerouting workflow.
    \item \textbf{Bridge Inventory} $\rightarrow$ \textit{``Can the vehicle physically traverse this alternate route?''} --- Enforces hard physical constraints (clearance, load limits) that no AI reasoning can override.
\end{itemize}

Together, these datasets answer the three fundamental middle-mile questions: \textit{risk}, \textit{trigger}, and \textit{feasibility}.

The two supporting datasets extend coverage: \textbf{DataCo} provides historical delay patterns for SRSNet forecasting model training, while the \textbf{Logistics Operations DB} (14 interconnected tables covering loads, trips, drivers, trucks, fuel, maintenance, and safety incidents) is vectorized into \texttt{SILVER.LOGISTICS\_VECTORIZED} to serve as the RAG knowledge base for natural language queries.

\subsection{Multimodal Ingestion Details}

\begin{table}[H]
\centering
\caption{Ingestion Methods by Dataset}
\small
\begin{tabularx}{\textwidth}{@{}l l l X@{}}
\toprule
\textbf{Dataset} & \textbf{Ingestion Method} & \textbf{Bronze Table} & \textbf{Format Details} \\
\midrule
US Accidents & \textbf{Snowpipe} (real-time) via S3 & \texttt{TRAFFIC\_INCIDENTS} & CSV $\rightarrow$ 47 columns (GPS, severity, road features, weather at scene) \\
NOAA GSOD & \textbf{External Table} (direct S3) & \texttt{NOAA\_GSOD} & Parquet $\rightarrow$ station-level daily summaries (PRCP, SNWD, VISIB, TEMP) \\
Bridge Inventory & \textbf{Snowpipe} via S3 & \texttt{BRIDGE\_INVENTORY} & CSV/GeoJSON $\rightarrow$ structural attributes + lat/lon coordinates \\
DataCo Supply Chain & \textbf{Snowpipe} via S3 & \texttt{RAW\_LOGISTICS} & CSV $\rightarrow$ 48 columns, Latin-1 encoding, normalized headers \\
Logistics Ops DB & \textbf{Snowpipe} via S3 & \texttt{LOGISTICS} & 14 CSVs $\rightarrow$ joined into single table with foreign-key resolution \\
\bottomrule
\end{tabularx}
\end{table}

\begin{itemize}
    \item \textbf{Snowpipe} provides event-driven, near-real-time ingestion for four datasets via S3 event notifications.
    \item \textbf{External Tables} avoid data duplication for the multi-terabyte NOAA archive --- queries run directly against S3.
    \item \textbf{Automated S3 Loading}: See \texttt{src/sql/02\_setup\_s3\_automation.sql} and \texttt{src/ingestion/setup\_s3\_automation.py}.
\end{itemize}

\subsection{Data Preprocessing and Sampling Strategy}

\begin{table}[H]
\centering
\caption{Preprocessing Scripts and Sampling Strategies}
\small
\begin{tabularx}{\textwidth}{@{}l l X X l@{}}
\toprule
\textbf{Dataset} & \textbf{Script} & \textbf{Key Transformations} & \textbf{Sampling Strategy} & \textbf{Silver Output} \\
\midrule
US Accidents & \texttt{preprocess\_accidents.py} & Aggregation by State/City/Hour; severity averaging; risk scoring & Geospatial clustering by proximity to freight corridors with severity weighting & \texttt{RISK\_HEATMAP\_VIEW} \\
\addlinespace
NOAA GSOD & \texttt{preprocess\_weather.py} & Rule-based alert extraction (PRCP $>$ 50mm, SNWD $>$ 10cm, VISIB $<$ 1mi); semantic description generation & Temporal windowing --- 4--8 hour adaptive patches aligned to middle-mile transit windows for SRSNet input & \texttt{WEATHER\_ALERTS} \\
\addlinespace
Bridge Inventory & \texttt{preprocess\_bridges.py} & Lat/lon $\rightarrow$ Snowflake \texttt{GEOGRAPHY} via \texttt{ST\_GEOGFROMTEXT}; constraint field extraction & \textbf{No sampling} --- full 600K+ inventory retained for exhaustive spatial join coverage (safety-critical) & \texttt{BRIDGE\_INVENTORY\_GEO} \\
\addlinespace
DataCo & \texttt{preprocess\_dataco.py} & Delay propensity scores; Z-score by shipping mode; late delivery risk cast & Full dataset retained for SRSNet training coverage & \texttt{CLEANED\_LOGISTICS} \\
\addlinespace
Logistics Ops & \texttt{preprocess\_logistics.py} & Record-to-text transformation; metadata JSON construction; chunk ID assignment & Full dataset --- all records vectorized for RAG retrieval completeness & \texttt{LOGISTICS\_VECTORIZED} \\
\bottomrule
\end{tabularx}
\end{table}

% ─────────────────────────────────────────────────────────────────────────────
\newpage
\section{Retrieval \& Processing Pipeline}
\label{app:retrieval}
% ─────────────────────────────────────────────────────────────────────────────

\noindent\textit{Full documentation: \texttt{docs/retrieval\_processing\_pipeline.md}}

\subsection{Chunking Strategy: Semantic Record-Level}

We employ a \textbf{Semantic Record-Level} chunking strategy rather than arbitrary text splitting. Structured database records from \texttt{BRONZE.LOGISTICS} are transformed into natural language sentences. Each ``chunk'' represents a complete logistics event (trip, maintenance record, or safety incident), preserving entity relationships.

\begin{itemize}
    \item \textbf{Strategy:} Record-to-Text Mapping
    \item \textbf{Implementation:} \texttt{src/preprocessing/preprocess\_logistics.py}
    \item \textbf{Rationale:} Logistics data is inherently transactional. A chunk is most meaningful when it represents a complete event, rather than a fixed-size text window.
\end{itemize}

\textbf{Example Transformation:}
\begin{lstlisting}[caption={Semantic chunking: structured data to natural language}]
# Input (structured):
Load_ID: 123, Origin: Chicago, Destination: St. Louis, Status: Late

# Output (chunk):
"Record type: load. Route: Chicago -> St. Louis, 300 miles.
 Status: Late. On-time: False."
\end{lstlisting}

Additional chunking strategies used across the pipeline:
\begin{itemize}
    \item \textbf{Time-Series:} Adaptive patching (4--8 hour windows) for SRSNet forecasting, aligned to middle-mile transit time scales.
    \item \textbf{Text:} 512--1024 token segments with sliding windows for longer document inputs.
    \item \textbf{Geospatial:} 10-mile route segments for localized accident risk scoring.
\end{itemize}

\subsection{Indexing and Retrieval Configuration}

The system uses a \textbf{Dense Vector Retrieval} configuration leveraging Snowflake's native vector support:

\begin{itemize}
    \item \textbf{Vector Type:} Snowflake \texttt{VECTOR(FLOAT, 768)} with L2 distance
    \item \textbf{Retrieval:} Cortex Search with ReMindRAG-guided traversal
    \item \textbf{Hybrid:} Dense/sparse retrieval with reranking for improved relevance
    \item \textbf{Metadata:} Chunks indexed alongside a JSON \texttt{METADATA} object containing \texttt{driver\_id}, \texttt{truck\_id}, \texttt{event\_date}, and \texttt{revenue} to support hybrid filtering.
    \item \textbf{Table:} \texttt{HYPERLOGISTICS\_DB.SILVER.LOGISTICS\_VECTORIZED}
\end{itemize}

\subsection{Embedding Models Used}

\begin{table}[H]
\centering
\caption{Embedding and LLM Models}
\begin{tabular}{@{}llrl@{}}
\toprule
\textbf{Model} & \textbf{Provider} & \textbf{Dimensions} & \textbf{Purpose} \\
\midrule
\texttt{snowflake-arctic-embed-m} & Snowflake Cortex & 768 & Primary embedding model for RAG vector search \\
\texttt{snowflake-arctic} & Snowflake Cortex & --- & LLM for grounded answer generation \\
Sentence Transformers & HuggingFace & 768 & Fallback for comparison/evaluation \\
\bottomrule
\end{tabular}
\end{table}

Running the embedding model within Snowflake eliminates the need to move data to external GPU clusters, ensuring low latency and high security.

\subsection{Example Retrieval Outputs}

\subsubsection{Query 1: ``What are the recent safety concerns for driver DRV\_102?''}

\begin{itemize}
    \item \textbf{Retrieved Chunk:} \textit{``Record type: safety. Safety incident: Hard Braking involving driver DRV\_102 on 2024-02-15. Details: Sudden stop on I-95 due to traffic surge.''}
    \item \textbf{Cortex AI Response:} \textit{``Based on the safety logs for Feb 15th, Driver DRV\_102 had a hard braking incident on I-95. The driver reported a traffic surge, and no damage was recorded. Recommendation: Schedule a quick safety refresh on maintaining distance in high-traffic zones.''}
\end{itemize}

\subsubsection{Query 2: ``Check if my route from Chicago to St.\ Louis is safe today.''}

\begin{itemize}
    \item \textbf{Retrieved Context:}
    \begin{itemize}
        \item \textit{Accident Data:} 12 high-severity incidents on I-55 in the last 24 hours.
        \item \textit{Weather Alert:} ``Station 724340 reports Heavy precipitation expected on 2024-02-27.''
    \end{itemize}
    \item \textbf{Cortex AI Response:} \textit{``Warning: The standard I-55 route from Chicago to St.\ Louis is currently high-risk. There are 12 active incident reports and heavy rain forecast for the next 6 hours. Recommendation: Delay departure by 4 hours or reroute via I-57 to avoid the storm cell.''}
\end{itemize}

\subsection{Preprocessing Scripts in GitHub}

All processing logic is automated via the following scripts in the repository:
\begin{itemize}
    \item \texttt{src/preprocessing/preprocess\_logistics.py} --- Semantic chunking \& vectorization prep
    \item \texttt{src/preprocessing/preprocess\_accidents.py} --- Risk heatmap aggregation
    \item \texttt{src/preprocessing/preprocess\_weather.py} --- Weather alert extraction
    \item \texttt{src/preprocessing/preprocess\_bridges.py} --- Spatial normalization (\texttt{GEOGRAPHY} conversion)
    \item \texttt{src/preprocessing/preprocess\_dataco.py} --- Delay propensity feature engineering
    \item \texttt{notebooks/srsnet\_training.ipynb} --- SRSNet model training notebook
\end{itemize}

% ─────────────────────────────────────────────────────────────────────────────
\newpage
\section{Application Integration}
\label{app:application}
% ─────────────────────────────────────────────────────────────────────────────

\noindent\textit{Full documentation: \texttt{docs/application\_integration.md}}

\subsection{Streamlit Project Interface}

The dashboard (\texttt{src/app/dashboard.py}) connects directly to Snowflake via Snowpark and provides four main tabs:

\begin{enumerate}
    \item \textbf{Risk Heatmap} --- Displays traffic incident hotspots across the US via Folium interactive maps, powered by \texttt{SILVER.RISK\_HEATMAP\_VIEW}.
    \item \textbf{Route Comparison} --- Side-by-side AI-suggested vs.\ original route with delay metrics and bridge clearance verification.
    \item \textbf{Evidence \& Reasoning} --- Shows the raw grounding data (accident risk, weather alerts, bridge constraints, logistics vectors) that the AI uses to form its answers.
    \item \textbf{Query Logs} --- Full audit trail of every AI query, response, grounding sources, and response time.
\end{enumerate}

The sidebar contains a natural language input (``Ask the AI Agent'') where users can submit logistics queries. The system:
\begin{enumerate}
    \item Retrieves relevant evidence from Silver-layer tables via \texttt{get\_evidence()}.
    \item Passes the evidence as context to \textbf{Snowflake Cortex} for grounded generation.
    \item Displays the grounded AI response with source citations.
    \item Logs the query, response, and metadata to \texttt{GOLD.QUERY\_LOGS}.
\end{enumerate}

\subsection{Evidence Display and Grounded Answer Generation}

When a user submits a query, the \texttt{get\_evidence()} function retrieves supporting data from three Silver-layer tables:

\begin{table}[H]
\centering
\caption{Evidence Sources for Grounded AI Responses}
\begin{tabularx}{\textwidth}{@{}l l X@{}}
\toprule
\textbf{Source Table} & \textbf{Data Type} & \textbf{Example Evidence} \\
\midrule
\texttt{RISK\_HEATMAP\_VIEW} & Accident risk aggregations & ``Miami, FL: 2,340 incidents, avg severity 2.8'' \\
\texttt{BRIDGE\_INVENTORY\_GEO} & Bridge compliance data & ``I-95 Bridge: 14.2ft clearance, 40 ton limit, Fair'' \\
\texttt{LOGISTICS\_VECTORIZED} & Operational text chunks & ``Trip T-5521: Chicago $\rightarrow$ St.\  Louis, 297 mi, on-time'' \\
\bottomrule
\end{tabularx}
\end{table}

This evidence is injected into the Cortex AI prompt as structured context, ensuring that the AI's answer is \textbf{grounded in real data} rather than hallucinated. The ``Evidence \& Reasoning'' tab (Tab 3) allows dispatchers to inspect the raw sources behind any recommendation.

\textbf{Example Grounded Response:}
\begin{quote}
\textbf{User Query:} ``What are the riskiest routes near Chicago?''

\textbf{AI Response:} ``Based on accident data, the Chicago metropolitan area has over 15,200 recorded traffic incidents with an average severity of 2.4. The Chicago to St.\ Louis corridor (I-55) shows the highest concentration of delays, with 3 late deliveries in the past month. Recommendation: Consider routing via I-65 through Indianapolis for time-sensitive shipments.''

\textit{Grounded on: SILVER.RISK\_HEATMAP\_VIEW, SILVER.LOGISTICS\_VECTORIZED}
\end{quote}

\subsection{Query Logging and Evaluation Outputs}

Every query submitted through the dashboard is automatically logged to \texttt{GOLD.QUERY\_LOGS}:

\begin{table}[H]
\centering
\caption{Query Log Schema}
\begin{tabular}{@{}ll@{}}
\toprule
\textbf{Column} & \textbf{Description} \\
\midrule
\texttt{QUERY\_ID}          & Unique identifier (UUID) \\
\texttt{USER\_ID}           & Snowflake username of the querying user \\
\texttt{QUERY\_TEXT}        & Raw natural language question \\
\texttt{RESPONSE\_TEXT}     & Full AI-generated response \\
\texttt{GROUNDING\_SOURCES} & JSON array of Silver tables used as evidence \\
\texttt{EXECUTION\_TIME\_MS} & End-to-end response time in milliseconds \\
\texttt{IS\_GROUNDED}       & Boolean --- TRUE if evidence was found and used \\
\texttt{CREATED\_AT}        & Timestamp of the query \\
\bottomrule
\end{tabular}
\end{table}

The ``Query Logs'' tab provides three summary metrics: \textbf{Total Queries}, \textbf{Grounded \%}, and \textbf{Avg Response Time}. These metrics support the evaluation framework defined in the Phase 1 proposal.

\subsection{Screenshots and Deployment}

\begin{itemize}
    \item \textbf{Local Deployment:} \texttt{streamlit run src/app/dashboard.py} $\rightarrow$ \texttt{http://localhost:8501}
    \item \textbf{Production Target:} Streamlit in Snowflake at \texttt{https://\textless account\textgreater.snowflakecomputing.com/streamlit/apps/HYPERLOGISTICS\_APP}
    \item \textbf{Technology Stack:} Streamlit + Folium (maps), Snowpark (backend), Cortex (AI engine), S3 $\rightarrow$ Snowpipe $\rightarrow$ Medallion (data pipeline)
\end{itemize}

% Screenshots can be added here if available:
% \begin{figure}[H]
%     \centering
%     \includegraphics[width=\textwidth]{screenshots/dashboard_risk_heatmap.png}
%     \caption{Risk Heatmap tab showing traffic incident hotspots}
% \end{figure}

% ─────────────────────────────────────────────────────────────────────────────
\newpage
\section{Snowflake Data Pipeline \& Schema}
\label{app:snowflake}
% ─────────────────────────────────────────────────────────────────────────────

\noindent\textit{Full documentation: \texttt{docs/snowflake\_data\_pipeline.md}}

\subsection{Snowflake Schema (Tables, Stages, Views)}

The database \texttt{HYPERLOGISTICS\_DB} follows a three-layer Medallion Architecture:

\begin{lstlisting}[caption={Medallion Architecture layout}]
HYPERLOGISTICS_DB
  +-- BRONZE (Raw Data)          -- Untouched data from S3
  +-- SILVER (Cleaned/Enriched)  -- Preprocessed, analytics-ready
  +-- GOLD   (Business Logic)    -- AI outputs, query logs, reports
\end{lstlisting}

\begin{table}[H]
\centering
\caption{Bronze Layer Tables}
\small
\begin{tabular}{@{}llrl@{}}
\toprule
\textbf{Table} & \textbf{Source Dataset} & \textbf{Records} & \textbf{Purpose} \\
\midrule
\texttt{TRAFFIC\_INCIDENTS} & US Accidents & $\sim$3.5M & GPS-tagged incidents for risk heatmaps \\
\texttt{BRIDGE\_INVENTORY} & National Bridge Inventory & 624K & Bridge clearance/load limits for safety veto \\
\texttt{RAW\_LOGISTICS} & DataCo Supply Chain & 180K & Historical delay patterns for SRSNet \\
\texttt{LOGISTICS} & Logistics Operations DB & 2.7M & Trucking ops for RAG knowledge base \\
\texttt{NOAA\_GSOD} (ext.) & NOAA GSOD & Millions & Weather disruption signals \\
\bottomrule
\end{tabular}
\end{table}

\begin{table}[H]
\centering
\caption{Silver Layer Tables}
\small
\begin{tabular}{@{}lll@{}}
\toprule
\textbf{Table/View} & \textbf{Source} & \textbf{Purpose} \\
\midrule
\texttt{RISK\_HEATMAP\_VIEW} & \texttt{TRAFFIC\_INCIDENTS} & Aggregated incident counts by State/City/Hour \\
\texttt{BRIDGE\_INVENTORY\_GEO} & \texttt{BRIDGE\_INVENTORY} & Bridge data with \texttt{GEOGRAPHY} type for spatial joins \\
\texttt{CLEANED\_LOGISTICS} & \texttt{RAW\_LOGISTICS} & Delay propensity scores for SRSNet training \\
\texttt{LOGISTICS\_VECTORIZED} & \texttt{LOGISTICS} & Text chunks + metadata for RAG retrieval \\
\texttt{WEATHER\_ALERTS} & \texttt{NOAA\_GSOD} & Semantic weather alerts (rain, snow, low visibility) \\
\bottomrule
\end{tabular}
\end{table}

\begin{table}[H]
\centering
\caption{Gold Layer Tables}
\begin{tabular}{@{}ll@{}}
\toprule
\textbf{Table} & \textbf{Purpose} \\
\midrule
\texttt{VECTOR\_EMBEDDINGS} & Cortex-generated 768-dim vectors for semantic search \\
\texttt{DISRUPTION\_REPORTS} & AI-generated disruption impact reports \\
\texttt{QUERY\_LOGS} & Full audit trail of every AI query \\
\bottomrule
\end{tabular}
\end{table}

\begin{table}[H]
\centering
\caption{Stages \& Pipes (Automated S3 Ingestion)}
\small
\begin{tabular}{@{}llll@{}}
\toprule
\textbf{Stage} & \textbf{S3 Path} & \textbf{Target Table} & \textbf{Pipe} \\
\midrule
\texttt{ACCIDENTS\_STAGE} & \texttt{s3://bucket/accidents/} & \texttt{TRAFFIC\_INCIDENTS} & \texttt{ACCIDENTS\_PIPE} \\
\texttt{BRIDGES\_STAGE} & \texttt{s3://bucket/bridges/} & \texttt{BRIDGE\_INVENTORY} & \texttt{BRIDGES\_PIPE} \\
\texttt{DATACO\_STAGE} & \texttt{s3://bucket/dataco/} & \texttt{RAW\_LOGISTICS} & \texttt{DATACO\_PIPE} \\
\texttt{LOGISTICS\_STAGE} & \texttt{s3://bucket/logistics/} & \texttt{LOGISTICS} & \texttt{LOGISTICS\_PIPE} \\
\texttt{NOAA\_S3\_STAGE} & \texttt{s3://noaa-gsod-pds/} & \texttt{NOAA\_GSOD} (ext.) & N/A \\
\bottomrule
\end{tabular}
\end{table}

\subsection{Reproducible Ingestion Scripts (SQL / Snowpark / Python)}

Scripts are run in the following order:

\begin{enumerate}
    \item \texttt{src/sql/00\_create\_database.sql} --- Creates database, 3 schemas, and all Bronze/Silver/Gold tables
    \item \texttt{src/sql/01\_setup\_noaa.sql} --- Creates external table pointing to public NOAA S3 bucket
    \item \texttt{src/sql/02\_setup\_s3\_automation.sql} --- Creates S3 stages and Snowpipe for automated ingestion
    \item \texttt{upload\_to\_s3.py} --- Uploads local CSV files to S3 (Snowpipe auto-loads)
    \item \texttt{src/preprocessing/preprocess\_*.py} --- Transforms Bronze $\rightarrow$ Silver via Snowpark Python
\end{enumerate}

\subsection{Example Queries Demonstrating Warehouse Connectivity}

\begin{lstlisting}[language=SQL, caption={Query 1: Top accident hotspots along freight corridors}]
SELECT STATE, CITY, INCIDENT_COUNT, AVG_SEVERITY
FROM HYPERLOGISTICS_DB.SILVER.RISK_HEATMAP_VIEW
ORDER BY INCIDENT_COUNT DESC LIMIT 10;
\end{lstlisting}

\begin{lstlisting}[language=SQL, caption={Query 2: Bridge clearance check for Illinois routes}]
SELECT FEATURES_DESC_006A AS BRIDGE_NAME,
       VERT_CLR_OVER_MT_053 AS VERTICAL_CLEARANCE_MT,
       OPERATING_RATING_064 AS LOAD_RATING, BRIDGE_CONDITION
FROM HYPERLOGISTICS_DB.SILVER.BRIDGE_INVENTORY_GEO
WHERE VERT_CLR_OVER_MT_053 < 4.5 AND STATE_CODE_001 = 17
ORDER BY VERT_CLR_OVER_MT_053 ASC;
\end{lstlisting}

\begin{lstlisting}[language=SQL, caption={Query 3: Delay analysis by shipping mode (Silver layer)}]
SELECT SHIPPING_MODE,
       AVG(DELAY_DAYS) AS avg_delay,
       AVG(DELAY_PROPENSITY) AS avg_propensity
FROM HYPERLOGISTICS_DB.SILVER.CLEANED_LOGISTICS
GROUP BY SHIPPING_MODE
ORDER BY avg_delay DESC;
\end{lstlisting}

\begin{lstlisting}[language=SQL, caption={Query 4: AI-powered query via Snowflake Cortex}]
SELECT SNOWFLAKE.CORTEX.COMPLETE('snowflake-arctic',
    'What are the main causes of supply chain delays
     in the Midwest?');
\end{lstlisting}

\begin{lstlisting}[language=SQL, caption={Query 5: Audit trail of AI interactions}]
SELECT QUERY_TEXT, RESPONSE_TEXT, GROUNDING_SOURCES,
       EXECUTION_TIME_MS, IS_GROUNDED, CREATED_AT
FROM HYPERLOGISTICS_DB.GOLD.QUERY_LOGS
ORDER BY CREATED_AT DESC LIMIT 10;
\end{lstlisting}

\subsection{How Snowflake Integrates with the Application}

\begin{table}[H]
\centering
\caption{Application Integration Points}
\small
\begin{tabularx}{\textwidth}{@{}l l X@{}}
\toprule
\textbf{Integration} & \textbf{Technology} & \textbf{Purpose} \\
\midrule
Data Ingestion & Snowpipe + S3 Events & Auto-loads CSV files from S3 into Bronze tables on arrival \\
Data Processing & Snowpark Python API & Preprocessing scripts transform data and write to Silver tables \\
AI Generation & Cortex (\texttt{COMPLETE}) & Dashboard queries the LLM for grounded answers \\
Vector Search & Cortex (\texttt{EMBED\_TEXT\_768}) & 768-dim embeddings for semantic retrieval in RAG \\
Dashboard & Streamlit + Snowpark & Real-time queries to Silver/Gold for visualizations and metrics \\
Observability & \texttt{GOLD.QUERY\_LOGS} & Every AI interaction logged with timing and grounding metadata \\
\bottomrule
\end{tabularx}
\end{table}

Snowflake serves as a \textbf{unified platform}: data storage, ML inference (via Cortex), and application hosting (via Streamlit in Snowflake) --- eliminating the need for external services or data movement.

% ─────────────────────────────────────────────────────────────────────────────
\newpage
\section{Reproducibility Plan}
\label{app:reproducibility}
% ─────────────────────────────────────────────────────────────────────────────

\noindent\textit{Full documentation: \texttt{docs/reproducibility\_plan.md} and \texttt{reproducibility/README\_repoductibility.md}}

\subsection{Environment Configuration}

\textbf{Requirements.txt} (pinned dependencies):
\begin{lstlisting}[caption={Key dependencies from requirements.txt}]
snowflake-snowpark-python>=1.0.0
streamlit>=1.0.0
torch>=2.0.0
transformers>=4.0.0
sentence-transformers>=2.0.0
pandas>=1.5.0
numpy>=1.24.0
boto3>=1.28.0
folium>=0.20.0
streamlit-folium>=0.26.0
python-dotenv>=1.0.0
\end{lstlisting}

\textbf{Setup Commands:}
\begin{lstlisting}[caption={Environment setup}]
git clone https://github.com/mosomo82/
    CS5542_SmartSC_Optimization_System.git
cd CS5542_SmartSC_Optimization_System
python -m venv venv
.\venv\Scripts\activate          # Windows
pip install -r requirements.txt
cp .env.template .env            # Fill in credentials
\end{lstlisting}

\textbf{External Service Requirements:}
\begin{itemize}
    \item \textbf{Snowflake} --- Enterprise Edition on AWS (Cortex AI, Snowpipe)
    \item \textbf{AWS S3} --- Standard tier for data staging
    \item \textbf{Python} --- 3.10 or 3.11 (Snowpark-compatible)
\end{itemize}

\textbf{Environment Variables:} Configured in \texttt{.env} (template at \texttt{.env.template}) with Snowflake credentials (\texttt{SNOWFLAKE\_ACCOUNT}, \texttt{SNOWFLAKE\_USER}, \texttt{SNOWFLAKE\_PASSWORD}, \texttt{SNOWFLAKE\_ROLE}, \texttt{SNOWFLAKE\_WAREHOUSE}, \texttt{SNOWFLAKE\_DATABASE}, \texttt{SNOWFLAKE\_SCHEMA}), AWS keys, S3 bucket name, and region.

\subsection{Model Versions and Dataset Versioning}

\begin{table}[H]
\centering
\caption{AI Model Versions (Snowflake-Managed --- No Local GPU Required)}
\begin{tabular}{@{}llll@{}}
\toprule
\textbf{Model} & \textbf{Provider} & \textbf{Version} & \textbf{Purpose} \\
\midrule
\texttt{snowflake-arctic} & Snowflake Cortex & Latest (managed) & LLM for grounded answers \\
\texttt{snowflake-arctic-embed-m} & Snowflake Cortex & Latest (managed) & 768-dim embeddings for RAG \\
\bottomrule
\end{tabular}
\end{table}

\begin{table}[H]
\centering
\caption{Dataset Versions}
\small
\begin{tabular}{@{}llrl@{}}
\toprule
\textbf{Dataset} & \textbf{Version} & \textbf{Records} & \textbf{Source} \\
\midrule
US Accidents & 2016--2023 (v14) & $\sim$7.7M & \href{https://www.kaggle.com/datasets/sobhanmoosavi/us-accidents}{Kaggle} \\
DataCo Supply Chain & v1 (2020) & $\sim$180K & \href{https://www.kaggle.com/datasets/shashwatwork/dataco-smart-supply-chain-for-big-data-analysis}{Kaggle} \\
National Bridge Inventory & NBI 2024 & $\sim$600K & \href{https://geodata.bts.gov/datasets/national-bridge-inventory/}{BTS/FHWA} \\
NOAA GSOD & Continuous (public) & Millions & \href{https://registry.opendata.aws/noaa-gsod/}{AWS Open Data} \\
Logistics Operations DB & v1 & $\sim$550K & \href{https://www.kaggle.com/datasets/yogape/logistics-operations-database}{Kaggle} \\
\bottomrule
\end{tabular}
\end{table}

\textbf{Data Versioning Strategy:}
\begin{itemize}
    \item \textbf{Bronze tables} contain raw, immutable snapshots of source CSVs.
    \item \textbf{Silver tables} are deterministically regenerated from Bronze by the preprocessing scripts.
    \item \textbf{Gold tables} contain runtime data (query logs) and are not expected to be identical across reproductions.
    \item SRSNet and ReMindRAG implementations are pinned to their NeurIPS 2025 paper versions.
\end{itemize}

\subsection{Random Seeds and Configuration Files}

The preprocessing pipeline is \textbf{deterministic by design} --- no random sampling or stochastic operations. All transformations use aggregations (\texttt{GROUP BY}, \texttt{AVG}, \texttt{COUNT}), column derivations (\texttt{withColumn}), and window functions (\texttt{Window.partitionBy}).

For SRSNet model training (\texttt{notebooks/srsnet\_training.ipynb}):
\begin{lstlisting}[language=Python, caption={Random seed configuration}]
import torch, numpy as np, random
SEED = 42  # Defined in src/config/config.yaml
random.seed(SEED); np.random.seed(SEED)
torch.manual_seed(SEED); torch.cuda.manual_seed_all(SEED)
torch.backends.cudnn.deterministic = True
\end{lstlisting}

\textbf{Configuration Files:}
\begin{table}[H]
\centering
\caption{Configuration Files}
\begin{tabular}{@{}ll@{}}
\toprule
\textbf{File} & \textbf{Purpose} \\
\midrule
\texttt{src/config/config.yaml} & Global seed (42), SRSNet hyperparams, ReMindRAG params, dataset paths \\
\texttt{.env} / \texttt{.env.template} & Snowflake and AWS credentials \\
\texttt{requirements.txt} & Pinned Python package versions \\
\texttt{.gitignore} & Ensures credentials and data files are not committed \\
\texttt{src/sql/00\_create\_database.sql} & Database schema definition (DDL) \\
\texttt{src/sql/01\_setup\_noaa.sql} & NOAA external table configuration \\
\texttt{src/sql/02\_setup\_s3\_automation.sql} & S3 stages, pipes, and integration setup \\
\bottomrule
\end{tabular}
\end{table}

\subsection{Run Instructions}

\textbf{Estimated Total Time: $\sim$2 hours}

\begin{enumerate}
    \item \textbf{Environment Setup ($\sim$10 min):} Clone repo, create venv, install \texttt{requirements.txt}, configure \texttt{.env} from template.

    \item \textbf{Snowflake Setup ($\sim$20 min):} Run SQL scripts \textit{in order} in Snowflake Worksheets:
    \begin{enumerate}[label=\alph*.]
        \item \texttt{src/sql/00\_create\_database.sql} --- Database, schemas, tables
        \item \texttt{src/sql/01\_setup\_noaa.sql} --- NOAA external table
        \item \texttt{src/sql/02\_setup\_s3\_automation.sql} --- S3 stages and Snowpipe
    \end{enumerate}

    \item \textbf{Data Loading ($\sim$30 min):} Download datasets from source links (Table 16), place CSVs in \texttt{data/}, then run \texttt{python upload\_to\_s3.py}.

    \item \textbf{Preprocessing ($\sim$15 min):}
\begin{lstlisting}
python src/preprocessing/preprocess_accidents.py
python src/preprocessing/preprocess_bridges.py
python src/preprocessing/preprocess_dataco.py
python src/preprocessing/preprocess_logistics.py
python src/preprocessing/preprocess_weather.py
\end{lstlisting}

    \item \textbf{Launch Dashboard ($\sim$5 min):} \texttt{streamlit run src/app/dashboard.py} $\rightarrow$ opens at \texttt{http://localhost:8501}

    \item \textbf{Verify:}
\begin{lstlisting}[language=SQL]
-- Check Bronze layer
SELECT COUNT(*) FROM HYPERLOGISTICS_DB.BRONZE.TRAFFIC_INCIDENTS;
-- Check Silver layer
SELECT COUNT(*) FROM HYPERLOGISTICS_DB.SILVER.RISK_HEATMAP_VIEW;
-- Test Cortex AI
SELECT SNOWFLAKE.CORTEX.COMPLETE('snowflake-arctic',
    'What causes supply chain delays?');
-- Check query logs
SELECT * FROM HYPERLOGISTICS_DB.GOLD.QUERY_LOGS
ORDER BY CREATED_AT DESC LIMIT 5;
\end{lstlisting}
\end{enumerate}

\begin{table}[H]
\centering
\caption{Expected Reproduction Checkpoints}
\begin{tabular}{@{}ll@{}}
\toprule
\textbf{Checkpoint} & \textbf{Expected Result} \\
\midrule
Database exists & \texttt{SHOW DATABASES LIKE 'HYPERLOGISTICS\%';} returns 1 row \\
Bronze tables populated & 4+ tables with data (3.5M accidents, 180K logistics, etc.) \\
Silver tables created & 5 tables/views with preprocessed data \\
Dashboard loads & Streamlit app at \texttt{localhost:8501} with 4 tabs \\
AI responds & Sidebar query returns a grounded response from Cortex \\
Query logs recorded & \texttt{GOLD.QUERY\_LOGS} shows logged queries \\
\bottomrule
\end{tabular}
\end{table}

% =============================================================================
% END OF APPENDICES — The \begin{thebibliography} section follows below
% =============================================================================


\newpage
\begin{thebibliography}{9}

\bibitem{srsnet}
Wu, X., Qiu, X., et al. (2025).
\textit{Enhancing Time Series Forecasting through Selective Representation Spaces: A Patch Perspective}.
arXiv preprint arXiv:2510.14510.
\url{https://arxiv.org/abs/2510.14510}

\bibitem{remindrag}
Hu, Y., Zhu, J., et al. (2025).
\textit{ReMindRAG: Low-Cost LLM-Guided Knowledge Graph Traversal for Efficient RAG}.
arXiv preprint arXiv:2510.13193.
\url{https://arxiv.org/abs/2510.13193}

\bibitem{cpp}
Maggiar, A., Dicker, L., \& Mahoney, M. W. (2024).
\textit{Consensus Planning with Primal, Dual, and Proximal Agents}.
Amazon Science.
\url{https://www.amazon.science/publications/consensus-planning-with-primal-dual-and-proximal-agents}

\bibitem{snowpark}
Snowflake Inc. (2025).
\textit{Snowpark Developer Guide: Python}.
\url{https://docs.snowflake.com/en/developer-guide/snowpark/python/index}

\bibitem{dataco}
Constante, F. (2019).
\textit{DataCo SMART SUPPLY CHAIN FOR BIG DATA ANALYSIS} [Dataset].
Mendeley Data.
\url{https://www.kaggle.com/datasets/shashwatwork/dataco-smart-supply-chain-for-big-data-analysis}

\bibitem{accidents}
Moosavi, S., Samavatian, M. H., Parthasarathy, S., \& Ramnath, R. (2019).
\textit{A Countrywide Traffic Accident Dataset}.
In Proceedings of the 27th ACM SIGSPATIAL.
\url{https://www.kaggle.com/datasets/sobhanmoosavi/us-accidents}

\bibitem{bridges}
Bureau of Transportation Statistics. (2025).
\textit{National Tunnel \& Bridge Inventory} [Dataset].
U.S. Department of Transportation.
\url{https://geodata.bts.gov/datasets/national-bridge-inventory/}

\bibitem{noaa}
National Oceanic and Atmospheric Administration. (2025).
\textit{Global Surface Summary of the Day (GSOD)} [Dataset].
AWS Open Data Registry.
\url{https://registry.opendata.aws/noaa-gsod/}

\bibitem{logistics}
Yogesh (2023).
\textit{Logistics Operations Database} [Dataset].
Kaggle.
\url{https://www.kaggle.com/datasets/yogape/logistics-operations-database}

\end{thebibliography}

\end{document}
